% !TeX program = pdflatex
% Topologia ogólna — samodzielny punkt wejścia rozdziału 16.
\documentclass[11pt,a4paper,oneside]{report}
\usepackage[utf8]{inputenc}
\usepackage[T1]{fontenc}
\usepackage{lmodern}
\usepackage[polish]{babel}
\usepackage[margin=2.5cm,headheight=14pt]{geometry}
\usepackage{amsmath,amssymb,amsthm,mathtools}
\usepackage{enumitem}
\usepackage{tikz}
\usetikzlibrary{arrows.meta,calc,positioning}
\usepackage{microtype}
\usepackage{placeins}
\usepackage[unicode]{hyperref}
\hypersetup{colorlinks=true,linkcolor=blue!45!black,urlcolor=blue!45!black}
\usepackage{fancyhdr}
\theoremstyle{plain}
\newtheorem{twierdzenie}{Twierdzenie}[section]
\newtheorem{lemat}[twierdzenie]{Lemat}
\newtheorem{stwierdzenie}[twierdzenie]{Stwierdzenie}
\newtheorem{wniosek}[twierdzenie]{Wniosek}
\theoremstyle{definition}
\newtheorem{definicja}[twierdzenie]{Definicja}
\newtheorem{przyklad}[twierdzenie]{Przykład}
\theoremstyle{remark}
\newtheorem{uwaga}[twierdzenie]{Uwaga}
\renewcommand{\proofname}{Dowód}
\newcommand{\R}{\mathbb{R}}
\newcommand{\Q}{\mathbb{Q}}
\setlength{\parindent}{0pt}
\setlength{\parskip}{4pt}
\setlength{\emergencystretch}{2em}
\pagestyle{fancy}
\fancyhf{}
\fancyhead[L]{\small Topologia ogólna}
\fancyhead[R]{\small\ifnum\value{chapter}>0 Rozdział \thechapter\else Spis treści\fi}
\fancyfoot[C]{\thepage}

% BEGIN GLOBAL CODE LISTINGS (generated from latex/pl/code-listings.tex)
\usepackage{xcolor}
\usepackage{listings}
\definecolor{codebg}{RGB}{248,249,251}
\definecolor{codeborder}{RGB}{180,191,203}
\definecolor{codecomment}{RGB}{61,113,79}
\definecolor{codekeyword}{RGB}{41,78,132}
\lstdefinelanguage{Julia}{
  morekeywords={abstract,baremodule,begin,break,catch,const,continue,do,else,elseif,
    end,export,false,finally,for,function,global,if,import,in,isa,let,local,macro,
    module,mutable,primitive,quote,return,struct,true,try,type,using,where,while,
    missing,nothing,println,print},
  sensitive=true,morecomment=[l]{\#},morecomment=[n]{\#=}{=\#},
  morestring=[b]",
}
\lstdefinestyle{skrypt}{
  backgroundcolor=\color{codebg},frame=single,rulecolor=\color{codeborder},
  framesep=5pt,frameround=tttt,framerule=.5pt,
  basicstyle=\ttfamily\footnotesize,keywordstyle=\color{codekeyword}\bfseries,
  commentstyle=\color{codecomment}\itshape,stringstyle=\color{purple!65!black},
  numbers=left,numberstyle=\tiny\color{gray!75!black},numbersep=8pt,
  stepnumber=1,breaklines=true,breakatwhitespace=false,
  columns=fullflexible,keepspaces=true,showstringspaces=false,tabsize=4,
  linewidth=.90\linewidth,captionpos=b,aboveskip=11pt,belowskip=12pt,
  xleftmargin=12pt,xrightmargin=2pt,
}
\lstdefinestyle{pythonpde}{style=skrypt,language=Python}
\lstdefinestyle{python}{style=skrypt,language=Python}
\lstdefinestyle{julia}{style=skrypt,language=Julia}
\lstset{style=skrypt}
\renewcommand{\lstlistingname}{Kod}
\renewcommand{\lstlistlistingname}{Spis kodów}
% END GLOBAL CODE LISTINGS
\begin{document}
\setcounter{chapter}{15}
\chapter{Funkcje ciągłe i wprowadzenie do uzwarceń}
\label{ch:topog-funkcje-uzwarcenia}

\noindent\textbf{Cel rozdziału.}
W rozdziale~5 zdefiniowaliśmy całkowitą regularność,
a w rozdziale~7 udowodniliśmy twierdzenie o zanurzaniu
przestrzeni Tychonoffa w kostkę $[0,1]^I$.
Nie powtarzamy tamtych dowodów. Wyjaśnimy teraz,
w jakim sensie topologia jest wyznaczona przez funkcje
ciągłe, czym jest uzwarcenie i kiedy dwa uzwarcenia
uznajemy za równoważne. Własności uniwersalne
konkretnych uzwarceń będą przedmiotem rozdziału~17.

\section{Topologia wyznaczona przez funkcje ciągłe}
\label{sec:topog-uzwarcenia-funkcje}

\begin{definicja}[Topologia początkowa rodziny funkcji]
\label{def:topog-uzwarcenia-topologia-poczatkowa}
Niech $X$ będzie zbiorem, a
$(f_i\colon X\to Y_i)_{i\in I}$ rodziną
odwzorowań do przestrzeni topologicznych $Y_i$.
\emph{Topologią początkową} tej rodziny
nazywamy najmniejszą topologię na $X$,
względem której wszystkie $f_i$ są ciągłe.
Tworzą ją jako podbazę zbiory
$f_i^{-1}(V)$, gdzie $i\in I$,
a $V$ jest otwarty w $Y_i$.
\end{definicja}
\begin{proof}
Przecięcia skończone zbiorów wskazanej
podbazy tworzą bazę pewnej topologii
(uwzględniamy puste przecięcie równe $X$).
W tej topologii każda $f_i$ jest ciągła
z definicji ciągłości przez przeciwobrazy.
Każda inna topologia zapewniająca
ciągłość wszystkich $f_i$ musi zawierać
wszystkie zbiory podbazy, a więc również
otwarte zbiory przez nie generowane.
\end{proof}

\begin{stwierdzenie}[Kryterium ciągłości]
\label{prop:topog-uzwarcenia-poczatkowa-ciaglosc}
Niech $\tau$ będzie topologią początkową
rodziny $(f_i\colon X\to Y_i)_{i\in I}$.
Dla dowolnej przestrzeni topologicznej $Z$
odwzorowanie $g\colon Z\to(X,\tau)$
jest ciągłe wtedy i tylko wtedy,
gdy wszystkie $f_i\circ g$ są ciągłe.
\end{stwierdzenie}
\begin{proof}
Jeśli $g$ jest ciągłe, to ciągłość złożeń
wynika z ciągłości wszystkich $f_i$.
Odwrotnie, przeciwobraz podbazowego zbioru
$f_i^{-1}(V)$ przez $g$ jest równy
$(f_i\circ g)^{-1}(V)$ i jest otwarty.
Przeciwobrazy skończonych przecięć
i dowolnych sum tych zbiorów też są otwarte.
Zatem $g$ jest ciągłe.
\end{proof}

\begin{twierdzenie}[Funkcyjny opis przestrzeni Tychonoffa]
\label{thm:topog-uzwarcenia-funkcyjny-opis}
Niech $X$ będzie przestrzenią topologiczną
spełniającą aksjomat $T_1$. Następujące
warunki są równoważne:
\begin{enumerate}[label=(\roman*)]
\item $X$ jest całkowicie regularna;
\item topologia $X$ jest topologią początkową
rodziny $C(X,[0,1])$ wszystkich ciągłych
funkcji $X\to[0,1]$;
\item funkcje z $C(X,[0,1])$ oddzielają punkty,
a dla każdego otwartego $U\subseteq X$
i każdego $x\in U$ istnieje $f\in C(X,[0,1])$
taka, że $f(x)=1$ i $f|_{X\setminus U}=0$.
\end{enumerate}
\end{twierdzenie}
\begin{proof}
Niech (i) zachodzi. Dla każdego
otwartego $U$ i $x\in U$ całkowita
regularność zastosowana do
$F=X\setminus U$ daje funkcję
$h\in C(X,[0,1])$ równą $0$ w $x$
i $1$ na $F$. Funkcja $f=1-h$
spełnia warunek z (iii).
Dla różnych $x,y$ zastosowanie tego
samego argumentu do $U=X\setminus\{y\}$
jest dozwolone przez $T_1$
i daje $f(x)=1$, $f(y)=0$.
Stąd (iii).

Warunek (iii) implikuje (ii).
Każda funkcja z $C(X,[0,1])$
jest ciągła w wyjściowej topologii,
zatem topologia początkowa jest
od niej nie silniejsza. Z drugiej strony
dla dowolnego otwartego $U$ i $x\in U$
wybieramy $f$ jak w (iii); wtedy
$x\in f^{-1}((1/2,1])\subseteq U$,
a zbiór pośrodku jest otwarty
w topologii początkowej.
Zatem obie topologie są identyczne.

Załóżmy (ii). Jeśli $F\subseteq X$
jest domknięty i $x\notin F$,
to $U=X\setminus F$ jest otwarty.
Z opisu przez podbazę istnieje
skończona rodzina funkcji
$f_1,\ldots,f_m\in C(X,[0,1])$
i otwarte $V_1,\ldots,V_m\subseteq[0,1]$
takie, że
\[
x\in W=\bigcap_{j=1}^{m}f_j^{-1}(V_j)
\subseteq U.
\]
Jeżeli $F=\varnothing$, funkcja stale
równa $0$ spełnia żądanie. W przeciwnym
razie $m\ge1$. Dla każdego $j$
istnieje ciągła funkcja
$h_j\colon[0,1]\to[0,1]$
równa $0$ w $f_j(x)$ i równa $1$
na $[0,1]\setminus V_j$;
jej istnienie wynika z całkowitej
regularności odcinka, wykazanej
w rozdziale~7.
Wówczas
$h=\max_{1\le j\le m}(h_j\circ f_j)$
jest ciągła, $h(x)=0$ i
$h|_{X\setminus W}=1$.
Ponieważ $F\subseteq X\setminus W$,
mamy $h|_F=1$. Jest to (i).
\end{proof}

\begin{przyklad}[Nie wystarcza samo oddzielanie punktów]
\label{ex:topog-uzwarcenia-samo-oddzielanie}
Niech $X=\R$ ma topologię dyskretną,
a w rodzinie funkcji pozostawmy
tylko jedno odwzorowanie
$f(x)=\frac12+\frac1\pi\arctan x$.
Funkcja ta jest ciągła i różnowartościowa,
więc oddziela punkty. Nie wyznacza jednak
topologii dyskretnej: jej topologia
początkowa jest zwykłą topologią
euklidesową na $\R$, ponieważ $f$
jest homeomorfizmem na $(0,1)$.
Rozdzielanie punktów przez
\emph{wybraną} rodzinę funkcji
nie gwarantuje więc odtworzenia
wszystkich zbiorów otwartych.
\end{przyklad}

\section{Definicja uzwarcenia i jego równoważności}
\label{sec:topog-uzwarcenia-definicje}

\begin{definicja}[Uzwarcenie]
\label{def:topog-uzwarcenia-uzwarcenie}
Niech $X$ będzie przestrzenią topologiczną.
\emph{Uzwarceniem Hausdorffa} przestrzeni $X$
nazywamy parę $(K,e)$, gdzie $K$ jest
zwartą przestrzenią Hausdorffa, a
$e\colon X\to K$ jest zanurzeniem
topologicznym o gęstym obrazie.
Przy utożsamieniu $X$ z $e(X)$
mówimy, że $X$ jest gęstą
podprzestrzenią $K$.
\end{definicja}

\begin{twierdzenie}[Istnienie uzwarcenia]
\label{thm:topog-uzwarcenia-istnienie}
Przestrzeń topologiczna $X$ ma uzwarcenie
Hausdorffa wtedy i tylko wtedy, gdy
jest przestrzenią Tychonoffa.
\end{twierdzenie}
\begin{proof}
Jeśli $(K,e)$ jest uzwarceniem,
to $K$ jest przestrzenią Tychonoffa:
zwarta przestrzeń Hausdorffa jest
normalna, a lemat Urysohna z rozdziału~5
zapewnia całkowitą regularność.
Zatem jej podprzestrzeń $e(X)$
jest Tychonoffa na mocy wyników
rozdziału~5, a więc także $X$.

Odwrotnie, dla przestrzeni Tychonoffa
stosujemy twierdzenie
o zanurzaniu w kostkę z rozdziału~7.
Daje ono zanurzenie
$e\colon X\to[0,1]^I$.
Domknięcie
$K=\overline{e(X)}^{\, [0,1]^I}$
jest zwarte na mocy twierdzenia
Tychonoffa i domkniętości
podzbioru zwartej przestrzeni;
jest również Hausdorffa,
a $e(X)$ jest w nim gęsty.
To jest dokładnie wniosek
o kostkowym uzwarceniu
z rozdziału~7.
\end{proof}

\begin{definicja}[Równoważność uzwarceń]
\label{def:topog-uzwarcenia-rownowaznosc}
Dwa uzwarcenia $(K_1,e_1)$
i $(K_2,e_2)$ tej samej przestrzeni $X$
są \emph{równoważne}, jeśli istnieje
homeomorfizm $H\colon K_1\to K_2$
spełniający $H\circ e_1=e_2$.
Wymóg zgodności z zanurzeniami
jest częścią definicji.
\end{definicja}

\begin{stwierdzenie}[Jednoznaczność homeomorfizmu]
\label{prop:topog-uzwarcenia-homeomorfizm-jednoznaczny}
Niech $(K_1,e_1)$ i $(K_2,e_2)$
będą uzwarceniami Hausdorffa przestrzeni $X$.
Istnieje co najwyżej jedno
odwzorowanie ciągłe $H\colon K_1\to K_2$
takie, że $H\circ e_1=e_2$.
W szczególności homeomorfizm
wykazujący ich równoważność
jest jednoznaczny.
\end{stwierdzenie}
\begin{proof}
Jeśli $H_1,H_2$ mają podaną własność,
to są równe na gęstym $e_1(X)$.
Zbiór punktów, w których
$H_1=H_2$, jest domknięty:
jest przeciwobrazem przekątnej
domkniętej w $K_2\times K_2$
(przestrzeń $K_2$ jest Hausdorffa).
Zawiera gęsty $e_1(X)$,
więc jest równy $K_1$.
\end{proof}

\begin{stwierdzenie}[Uzwarcenie przestrzeni zwartej]
\label{prop:topog-uzwarcenia-zwarta}
Jeśli $X$ jest zwartą przestrzenią
Hausdorffa, to każde jej uzwarcenie
Hausdorffa $(K,e)$ spełnia $e(X)=K$.
Tym samym każde takie uzwarcenie
jest równoważne $(X,\mathrm{id}_X)$.
\end{stwierdzenie}
\begin{proof}
Obraz $e(X)$ jest zwarty jako obraz
ciągły zwartej przestrzeni,
a więc domknięty w przestrzeni
Hausdorffa $K$. Jest zarazem gęsty,
więc $e(X)=K$. Zanurzenie
$e\colon X\to K$ jest zatem
homeomorfizmem i zapewnia równoważność.
\end{proof}

\section{Porządek uzwarceń}
\label{sec:topog-uzwarcenia-porzadek}

\begin{definicja}[Porównywanie uzwarceń]
\label{def:topog-uzwarcenia-porzadek}
Niech $(K_1,e_1)$ i $(K_2,e_2)$ będą
uzwarceniami Hausdorffa przestrzeni $X$.
Mówimy, że $(K_1,e_1)$ jest
\emph{nie mniejsze} od $(K_2,e_2)$,
i piszemy $(K_2,e_2)\preccurlyeq(K_1,e_1)$,
jeżeli istnieje ciągłe
$p\colon K_1\to K_2$ spełniające
$p\circ e_1=e_2$.
\end{definicja}

\begin{stwierdzenie}[Porządek na klasach uzwarceń]
\label{prop:topog-uzwarcenia-porzadek-klasy}
Niech $X$ będzie przestrzenią Tychonoffa.
Jeśli $(K_2,e_2)\preccurlyeq(K_1,e_1)$,
to odwzorowanie $p\colon K_1\to K_2$
spełniające $p\circ e_1=e_2$
jest wyznaczone jednoznacznie
i jest surjektywne.
Ponadto relacja $\preccurlyeq$
jest zwrotna i przechodnia;
jeżeli zachodzi ona w obu kierunkach,
to uzwarcenia są równoważne.
Wobec tego wyznacza porządek częściowy
na klasach równoważności uzwarceń.
\end{stwierdzenie}
\begin{proof}
Jednoznaczność $p$ wynika ze
stwierdzenia o jednoznaczności
odwzorowań ciągłych zgodnych
z zanurzeniami z poprzedniej sekcji.
Zbiór $p(K_1)$ jest zwarty jako obraz
ciągły zwartego $K_1$, a więc
domknięty w przestrzeni Hausdorffa $K_2$.
Zawiera $p(e_1(X))=e_2(X)$,
czyli gęsty podzbiór $K_2$.
Stąd $p(K_1)=K_2$.

Zwrotność wynika z możliwości
wyboru $p=\mathrm{id}_{K}$.
Jeżeli $p\colon K_1\to K_2$
oraz $q\colon K_2\to K_3$
są ciągłe i spełniają odpowiednio
$p\circ e_1=e_2$ i $q\circ e_2=e_3$,
to $(q\circ p)\circ e_1=e_3$,
co dowodzi przechodniości.

Załóżmy teraz istnienie
ciągłych $p\colon K_1\to K_2$
i $q\colon K_2\to K_1$ zgodnych
z zanurzeniami w obie strony.
Wtedy
$(q\circ p)\circ e_1=e_1$
oraz $(p\circ q)\circ e_2=e_2$.
Oba złożenia są zgodne z identycznościami
na gęstych obrazach $e_1(X)$ i $e_2(X)$.
Z wcześniej udowodnionej jednoznaczności
wynika $q\circ p=\mathrm{id}_{K_1}$
i $p\circ q=\mathrm{id}_{K_2}$.
Zatem $p$ jest homeomorfizmem,
a uzwarcenia są równoważne.

Na klasach równoważności definicja
nie zależy od wyboru reprezentantów:
przy zamianie uzwarceń na równoważne
wystarczy złożyć $p$ z odpowiednimi
homeomorfizmami zgodnymi z zanurzeniami.
Dowiedzione zwrotność, przechodniość
i antysymetria na tych klasach
dają porządek częściowy.
\end{proof}

\section{Uzwarcenia nierównoważne: przykład}
\label{sec:topog-uzwarcenia-przyklady}

\begin{przyklad}[Dwa uzwarcenia przedziału otwartego]
\label{ex:topog-uzwarcenia-dwa-przedzialu}
Niech $X=(0,1)$ z topologią euklidesową.
Pierwszym uzwarceniem jest
$K_1=[0,1]$ z zanurzeniem
$e_1(t)=t$.

Drugim jest okrąg jednostkowy
$K_2=S^1=\{z\in\mathbb C:|z|=1\}$
z zanurzeniem
$e_2(t)=\exp(2\pi it)$.
Jest ono różnowartościowe na $(0,1)$:
jeśli $e_2(s)=e_2(t)$, to $s-t$
jest liczbą całkowitą z $(-1,1)$,
zatem $s=t$.
Jest homeomorfizmem na obraz
$S^1\setminus\{1\}$:
odwrotnością jest ciągły argument
o wartościach w $(0,2\pi)$,
podzielony przez $2\pi$.
Obraz jest gęsty w $S^1$,
więc $(K_2,e_2)$ jest uzwarceniem.

Te uzwarcenia nie są równoważne.
Gdyby homeomorfizm
$H\colon[0,1]\to S^1$ spełniał
$H(t)=e^{2\pi it}$ dla $t\in(0,1)$,
to z ciągłości w $0$ i $1$
mielibyśmy $H(0)=H(1)=1$.
Jest to sprzeczne z różnowartościowością
homeomorfizmu.
W ten sposób końce przedziału
mogą zostać dodane jako dwa
różne punkty albo jako jeden punkt.
\end{przyklad}

\begin{uwaga}[Zakres kolejnego rozdziału]
W rozdziale~17 zbadamy dokładniej
uzwarcenie Aleksandrowa oraz
uzwarcenie Čecha--Stone'a, w tym ich
założenia i własności uniwersalne.
Nie identyfikujemy na tym etapie
dowolnego uzwarcenia z tymi konstrukcjami:
przykład powyżej pokazuje, że
same zwartość i gęstość obrazu
nie wyznaczają uzwarcenia jednoznacznie.
\end{uwaga}

\end{document}
